\documentclass[12pt]{article}
\usepackage[T1]{fontenc}
\usepackage[utf8]{inputenc}
\usepackage{lmodern}
\usepackage{textcomp}
\usepackage{enumitem}
\usepackage{geometry}
\geometry{margin=1in}
\begin{document}
\begin{center}
Answer Key - Biology - Graduate Level Assessment 23 \
\small (Answers are ordered exactly as in the question paper.)
\end{center}

\section*{Multiple Choice}
\begin{enumerate}[label=\arabic*.]
\item \textbf{Answer:} A
\\ \textit{Options:} A) peroxisome; B) cytoplasm; C) nucleus; D) intermembrane space; E) outer mitochondrial membrane
\\ \textit{Note:} The correct answer is not accurate, as the Krebs cycle, also known as the citric acid cycle, actually occurs in the mitochondria of human cells, where it plays a crucial role in cellular respiration and energy production.
\item \textbf{Answer:} A
\\ \textit{Options:} A) polyspermy reaction; B) enzymatic reaction; C) proteolytic reaction; D) acrosomal reaction; E) oocyte reaction
\\ \textit{Note:} The correct answer is the polyspermy reaction because the rise in intracellular free calcium triggers a series of events, including the release of proteolytic enzymes, that function to modify the oocyte's extracellular environment and prevent multiple sperm from fertilizing the egg, thereby ensuring successful fertilization.
\item \textbf{Answer:} A
\\ \textit{Options:} A) positive chemotropism; B) positive phototropism; C) negative phototropism; D) negative chemotropism; E) positive hydrotropism
\\ \textit{Note:} The plant's growth in the opposite direction of gravitational force is an example of positive chemotropism because it reflects the plant's ability to grow towards favorable chemical gradients, which can be influenced by gravitational signals.
\item \textbf{Answer:} A
\\ \textit{Options:} A) parallel selection.; B) sympatric selection.; C) artificial selection.; D) allopatric selection.; E) directional selection.
\\ \textit{Note:} Parallel selection leads to sexual dimorphism as it results from similar selective pressures acting on both sexes, often in response to environmental factors or mating strategies, thereby promoting distinct traits that enhance reproductive success.
\item \textbf{Answer:} B
\\ \textit{Options:} A) Flowers with a narrow opening that only small insects can access; B) Pendant (hanging) red-colored flowers; C) Modified petals to provide a landing space; D) Flowers with a strong sweet fragrance; E) Bright blue flowers that open in the early morning
\\ \textit{Note:} Pendant red-colored flowers are adapted to attract hummingbirds, which are drawn to bright colors and have the ability to access nectar from flowers that are structured for their feeding style, whereas bees are less responsive to red hues and prefer flowers that are more accessible and brightly colored in the blue and yellow spectrum.
\end{enumerate}

\section*{True / False}
\begin{enumerate}[label=\arabic*.]
\setcounter{enumi}{5}
\item \textbf{Answer:} False
\\ \textit{Options:} A) True; B) False
\\ \textit{Note:} Increasing the treatment volume to include a longer nerve length for trigeminal neuralgia radiosurgery does not significantly improve pain relief but may increase complications.
\item \textbf{Answer:} True
\\ \textit{Options:} A) True; B) False
\\ \textit{Note:} This study showed for the first time that RALP induces lower tissue trauma than RRP.
\item \textbf{Answer:} True
\\ \textit{Options:} A) True; B) False
\\ \textit{Note:} We conclude that GES is more effective in improving long-term GI symptoms and costs, and decreasing use of healthcare resources than intensive medical therapy, in this sample of patients with the symptoms of GP followed for 3 years. Certain patients with GP form a high-risk group in terms of costs, quality of life, morbidity and mortality.
\item \textbf{Answer:} True
\\ \textit{Options:} A) True; B) False
\\ \textit{Note:} MRI of children with OCD consistently showed secondary physis disruption, overlying chondroepiphysial widening, and subchondral bone edema. We suggest that disruption of normal endochondral ossification may be associated with juvenile OCD.
\item \textbf{Answer:} False
\\ \textit{Options:} A) True; B) False
\\ \textit{Note:} In this study, we found no systematic association between isolated episodes of binge drinking during early pregnancy and child motor function at age 5.
\end{enumerate}

\section*{Long Form Response}
\begin{enumerate}[label=\arabic*.]
\setcounter{enumi}{10}
\item \textbf{Answer:} Red-green color blindness is an X-linked recessive trait, meaning that males (XY) are affected if they inherit the mutant allele on their single X chromosome, while females (XX) require two copies to express the condition. In this scenario, the red-green color- blind man ($X^cY$) has a normal vision partner who is a carrier for the trait ($X^NX$^c), as her father was color blind. The potential offspring would exhibit the following phenotypic ratios: 1/4 of the offspring will be females with normal vision, 1/4 will be males with normal vision, 1/4 will be females who are color blind, and 1/4 will be males who are color blind. This results from the combination of maternal X chromosomes and the paternal Y chromosome, illustrating the inheritance patterns of X-linked traits.
\\ \textit{Note:} correct because it accurately identifies the X-linked recessive inheritance pattern of red-green color blindness and correctly determines the genotypes of the parents, leading to the expected phenotypic ratios of the offspring, including the likelihood of color-blindness in males and carrier status in females.
\item \textbf{Answer:} Algae are classified across the kingdoms Fungi, Plantae, and Archaea based on their cellular structure, reproductive methods, and ecological roles. In the Plantae kingdom, algae are primarily photoautotrophic organisms, characterized by their ability to perform photosynthesis, with examples including green algae (Chlorophyta) and red algae (Rhodophyta). Conversely, some algae, such as the lichens formed by cyanobacteria and fungi, align more closely with the Fungi kingdom due to their symbiotic relationships and heterotrophic characteristics. Furthermore, certain archaea exhibit photosynthetic capabilities, suggesting a broader understanding of photosynthetic processes that encompasses the Archaea domain. These classifications underscore the evolutionary relationships and functional diversity of algae, highlighting their ecological importance and the necessity for a nuanced view of organismal classification beyond traditional taxonomic boundaries.
\\ \textit{Note:} correct because it accurately identifies the classification of algae based on their photosynthetic capabilities, cellular characteristics, and ecological relationships, which determine their placement in the Plantae kingdom and their associations with Fungi through lichens.
\end{enumerate}

\vfill
\noindent\textit{End of Answer Key}
\end{document}